% !TeX root = ../深度学习笔记.tex
% 由主文档通过 \input 载入；请编译主文档。
\chapter*{前言与全书路线}
\addcontentsline{toc}{chapter}{前言与全书路线}

\section*{课程定位与学习目标}
本书面向具有一元微积分、初步线性代数和基础编程知识的深度学习初学者，目标是建立从问题建模、网络设计、梯度计算到实验验证的完整认识。\emph{深度学习（deep learning）}以多层可学习变换构造数据表示，既涉及函数的表达能力，也涉及从有限数据中估计参数的方法。学习重点不是记忆网络名称，而是能够回答：输入和输出是什么，模型施加了什么结构假设，目标函数衡量什么，参数如何更新，以及实验支持多大范围的结论。

本书按问题、数学工具、网络训练、典型架构和进阶应用组织正文。基础部分以可手算的小模型建立直觉，关键结论给出假设与必要推导；核心部分强调数学表达与程序计算的一致性，进阶部分介绍生成模型、结构数据和可靠性问题。各章配有解答完整的例题及概念、推导和实验练习，实验以明确的输入、对照和评价协议为任务要求，不预设尚未运行的结果。紧凑性体现在压缩重复叙述，保留定义条件、推导依据和算法边界。

\section*{六篇二十四章的课程总纲}
\begin{enumerate}
  \item \textbf{第一篇：问题与数学基础。}第\ref{chap:dl-introduction}章“深度学习的问题、方法与学习路线”；第\ref{chap:dl-mathematics}章“线性代数、微积分与数值计算”；第\ref{chap:dl-probability}章“概率、统计与信息量”；第\ref{chap:dl-baselines}章“监督学习与线性基线”。从具体任务出发，建立贯穿全书的模型、数据、损失与评估语言。
  \item \textbf{第二篇：神经网络与训练原理。}第\ref{chap:dl-mlp}章“前馈神经网络与表示能力”；第\ref{chap:dl-backprop}章“计算图、反向传播与自动微分”；第\ref{chap:dl-optimization}章“优化算法与训练动力学”；第\ref{chap:dl-regularization}章“初始化、归一化与正则化”；第\ref{chap:dl-engineering}章“训练程序、实验设计与故障诊断”。完成从手工求导到可靠训练一个网络的闭环。
  \item \textbf{第三篇：典型网络架构。}第\ref{chap:dl-cnn}章“卷积神经网络”；第\ref{chap:dl-vision}章“残差网络与视觉任务”；第\ref{chap:dl-rnn}章“循环神经网络与序列建模”；第\ref{chap:dl-attention}章“注意力机制”；第\ref{chap:dl-transformer}章“Transformer 与序列架构”。比较局部连接、状态递推和内容寻址三种信息组织方式。
  \item \textbf{第四篇：表示学习、预训练与应用。}第\ref{chap:dl-representation}章“表示学习与自监督学习”；第\ref{chap:dl-transfer}章“迁移学习、预训练与微调”；第\ref{chap:dl-language}章“语言建模与生成”；第\ref{chap:dl-multimodal}章“多模态学习与跨模态对齐”。把网络结构与学习任务分开，解释数据、目标函数和适配方式如何共同决定模型用途。
  \item \textbf{第五篇：深度生成模型。}第\ref{chap:dl-vae}章“概率生成建模与变分自编码器”；第\ref{chap:dl-gan-flow}章“生成对抗网络与正规化流”；第\ref{chap:dl-diffusion}章“扩散模型与去噪学习”。沿概率建模、训练目标、采样过程和评价方式比较不同生成路线。
  \item \textbf{第六篇：进阶专题与综合实践。}第\ref{chap:dl-gnn}章“图神经网络与结构化数据”；第\ref{chap:dl-reliability}章“泛化、鲁棒性与不确定性”；第\ref{chap:dl-projects}章“效率、部署与综合项目”。把模型放回计算预算、分布变化和真实使用条件中检验。
\end{enumerate}

\section*{阅读路线与贯穿案例}
基础路线依次学习第\ref{chap:dl-introduction}至第\ref{chap:dl-engineering}章，再完成卷积网络及视觉任务；序列路线在训练基础上学习第\ref{chap:dl-rnn}至第\ref{chap:dl-transformer}章，再进入表示学习、迁移学习和语言建模。进阶读者可在概率基础、反向传播与训练程序熟练后进入生成模型；图神经网络是可独立选择的结构数据专题。第\ref{chap:dl-reliability}章中有关泛化与评估的内容应在每次实验时回看，第\ref{chap:dl-projects}章用于整合全书方法。

贯穿案例分成三条：二维合成数据用于观察决策边界和核对梯度，小型图像分类用于比较网络结构与正则化，短序列预测用于区分训练目标和生成行为。生成模型先使用二维分布再使用小图像，避免把视觉效果当成数学解释。实验先建立简单基线，再一次改变一个主要因素；记录数据划分、参数量、训练预算、随机种子、评价指标与失败样本。基础阅读可结合文献\cite{goodfellow2016}，实现练习可结合文献\cite{d2l}，专门模型另列原始论文。

\section*{全书记号与论证约定}
标量用普通字母表示，向量与矩阵用粗体表示，单个样本的特征向量默认是列向量。训练集记为 $\mathcal D=\{(\mathbf x_i,y_i)\}_{i=1}^{n}$，其中 $n$ 为样本数，$\mathbf x_i\in\mathbb R^d$，$d$ 为输入维度；目标 $y_i$ 的取值空间随任务定义。随机输入、随机目标分别记为 $\mathbf X,Y$，其联合分布记为 $P$。批量数据矩阵另记为 $\mathbf X_{\mathrm{batch}}\in\mathbb R^{B\times d}$，每行是一个样本，$B$ 为批量大小；不要把这个矩阵与随机向量 $\mathbf X$ 混同。

模型写为 $f_\theta$，$\theta\in\mathbb R^p$ 汇集全部可训练参数，$p$ 为参数总数；样本损失写为 $\ell(f_\theta(\mathbf x),y)$，训练目标写为 $J(\theta)$。参数、超参数和优化器状态分别定义。层索引用 $l$，迭代或时间索引用 $t$，其范围在各章说明；$\mathbf I_d$ 表示 $d$ 阶单位矩阵，$\|\cdot\|_2$ 表示欧氏范数，$\odot$ 表示逐元素乘法，$\log$ 默认是自然对数。期望、梯度和概率式都注明所针对的随机量或变量。理论结论注明条件，经验规律通过对照实验讨论，算法输出与全局最优解不作无条件等同。

